\chapter{What we are building}
\label{ch:building}

We are building \texttt{arbiter} (ARBITER: Agent Runtime for Bounded Inference, Triage,
Evaluation \& Routing): an open-source Python SDK that adds \emph{typed, calibrated
decision points} to an LLM agent a developer already runs (LangChain, Pydantic AI, or the
OpenAI Agents SDK).

\section{The problem in one paragraph}

An LLM agent makes many small, bounded decisions on every loop: which tool to call next,
whether the task is finished, whether it is stuck. Today the same expensive model that
plans and writes text also makes these decisions. That is slow, costly, and, worst of all,
the decision comes back as plain text with \emph{no usable number} telling you how sure the
model is. You cannot tell a confident choice from a coin flip.

\section{The idea}

At each decision point the SDK does four things (Figure~\ref{fig:value}):

\begin{enumerate}[itemsep=2pt]
  \item Build a small, token-budgeted \term{projection} of the agent's state.
  \item Ask a cheap, swappable \term{engine} (Jev by default) a bounded typed question
        and get back \emph{probabilities}.
  \item A deterministic \term{policy} reads the (recalibrated) probability and either
        \textbf{acts}, \textbf{narrows} the LLM's options, or \textbf{escalates} to the
        LLM or a human.
  \item Log the decision and its outcome so thresholds can be fitted and results reproduced.
\end{enumerate}

\begin{figure}[t]\centering
\begin{tikzpicture}[node distance=8mm and 16mm]
  \node[frontend] (state) {agent state};
  \node[infra, right=of state] (proj) {projection\\(\,$\le$32k tokens)};
  \node[model, right=of proj] (engine) {engine\\(Jev default)};
  \node[backend, right=of engine] (policy) {policy};
  \node[frontend, above right=6mm and 12mm of policy] (act) {act};
  \node[frontend, right=of policy] (narrow) {narrow options};
  \node[security, below right=6mm and 12mm of policy] (esc) {escalate\\(LLM / human)};
  \draw[flow] (state) -- (proj);
  \draw[flow] (proj) -- (engine);
  \draw[flow] (engine) -- node[flowlabel]{prob.} (policy);
  \draw[flow] (policy) -- (act);
  \draw[flow] (policy) -- (narrow);
  \draw[flow] (policy) -- (esc);
\end{tikzpicture}
\caption{The arbiter value loop: state becomes a small projection, a cheap engine
returns probabilities, and a deterministic policy acts, narrows, or escalates.\datalabel{Illustrative}}
\label{fig:value}
\end{figure}

\section{Why this and not something else}

The research phase (Part~IV) reached a clear verdict: integrations already exist
(LangChain, Pydantic AI, and Vercel all ship Jev support, and about 200 community projects
exist~\citep{langchain_typesafe,pydantic_typesafe}). What nobody has published is
\emph{controlled evidence}, with proper baselines and trajectory-level calibration, that
inserting a calibrated decision component actually lowers cost per successful task without
hurting reliability. So the project has two co-equal deliverables:

\begin{description}[leftmargin=1.4cm,style=nextline]
  \item[Deliverable A --- the SDK] the instrument developers install.
  \item[Deliverable B --- the evidence] a pre-registered, reproducible study (with a
    released decision dataset and trajectory-checkpoint calibration) that no one else has.
\end{description}

\begin{whyhere}
The engine is \emph{swappable} on purpose. If Jev turns out not to help on your workload,
you change one line to an LLM-backed engine, a small trained classifier, or the
open-weights Laya model~\citep{laya}, and the SDK still works. The product survives a
negative result about any single engine.
\end{whyhere}

\section{What it is not}

It is not a new agent framework (that is an explicit non-goal). It does not plan, write
tool arguments, or generate text; the LLM keeps those jobs. It is never the security or
authorization boundary; deterministic code owns that. The MVP is deliberately small: two
decision points (which tool, and done/stuck) on one host framework.

\begin{checkbox}
Why can a model's answer be right while its confidence is untrustworthy? What would you
measure to catch that? (Answer in Appendix~\ref{app:answers}.)
\end{checkbox}
